\documentclass[11pt, a4paper, roman]{moderncv}

% --- ESTILO DE ECONOMIA DE ESPAÇO ---
% O estilo 'banking' remove a coluna lateral de datas e usa a largura total
\moderncvstyle{banking}
\moderncvcolor{blue}

\usepackage[utf8]{inputenc}
\usepackage[T1]{fontenc}
\usepackage{geometry}
\usepackage{textcomp}
\usepackage{newtxtext} % Times New Roman
\usepackage{enumitem} % Para melhor controle das listas
\usepackage{ragged2e}
\usepackage{needspace}
\hyphenpenalty=10000 \exhyphenpenalty=10000 % sem hifenização

% Listas: alinhadas à esquerda (sem hifenização) e com o mesmo espaçamento em todas as entradas
\setlist[itemize]{leftmargin=1.3em, topsep=0.2em, itemsep=0.15em, parsep=0pt, partopsep=0pt, before=\RaggedRight}

% --- AJUSTE DAS MARGENS (Máximo aproveitamento) ---
\geometry{left=0.5in, right=0.5in, top=0.5in, bottom=0.5in}

% --- DADOS PESSOAIS ---
\firstname{Orlando}
\familyname{Mota Pires}

\email{orlandomota2002@gmail.com}
\social[linkedin]{orlando-mota-0360a5209}
\social[github]{orlandomotapires}
\homepage{https://orlandomotapires.github.io}

% Links sublinhados, como nas versões anteriores
\newcommand{\cvlink}[2]{\underline{\href{#1}{#2}}}

% Entrada com o cargo primeiro: {cargo}{datas}{instituição}{local}{detalhes}
% Linha 1: cargo em negrito, datas à direita. Linha 2: instituição em itálico, local à direita.
\newcommand{\cvrole}[5]{\needspace{5\baselineskip}\cventry[0.85em]{#4}{#3}{#1}{\mdseries #2}{}{#5}}

% Publicação: {título}{autores}{veículo e links}
\setlength{\listitemsymbolspace}{0.3em}
\newcommand{\pub}[3]{\cvlistitem[0.7em]{\RaggedRight\textbf{#1}\\ #2\\ #3}}

% --- INÍCIO DO DOCUMENTO ---
\begin{document}

% --- CABEÇALHO ---
\makecvtitle
\vspace*{-3.0em} % Reduz o espaço logo após o título
\begin{center}
Computer engineer working on hybrid quantum--classical algorithms, high-performance computing and biosignals.\\
Looking for a research position in quantum computing and a master's programme.\\[0.4em]
\textbf{Portuguese}: Native \quad | \quad \textbf{English}: C1 \quad | \quad \textbf{German}: B2
\end{center}

% --- SEÇÃO: FORMAÇÃO ACADÊMICA ---
\section{Education}

\cvrole
  {Bachelor's Degree in Computer Engineering}
  {Feb 2021 -- Dec 2026}
  {Universidade Senai Cimatec}
  {Salvador, Brazil}
  {
  \begin{itemize}
    \item GPA 8.58/10.00
    \item Thesis: \cvlink{https://orlandomotapires.github.io/assets/pdf/Hybrid_Quantum_Classical_Graph_Partitioning_Algorithm_for_Domain_Decomposition.pdf}{Hybrid Quantum-Classical Graph Partitioning Algorithm for Domain Decomposition}, carried out at Fraunhofer IPT. Supervised by MSc.\ Stefan Schröder and Dr.\ Anderson Buarque.
  \end{itemize}
  }

\cvrole
  {Exchange Student}
  {Dec 2024 -- Aug 2026}
  {RWTH Aachen University}
  {Aachen, Germany}
  {}

\cvrole
  {JUNIQ-EPIQ Summer School on Quantum Computing}
  {Sep 2025}
  {Jülich Supercomputing Centre (JSC)}
  {Jülich, Germany}
  {
  \begin{itemize}
    \item Solved \cvlink{https://orlandomotapires.github.io/assets/pdf/Teilnahmebest\%C3\%A4tigung\%20Mota\%20Pires.pdf}{max-cut and travelling salesman problems} on gate-based quantum devices and quantum annealers.
  \end{itemize}
  }

% --- SEÇÃO: EXPERIÊNCIAS ---
\section{Research Experience}

\cvrole
  {Student Assistant}
  {Feb 2025 -- Sep 2026}
  {Fraunhofer Institute for Production Technology IPT}
  {Aachen, Germany}
  {
  \begin{itemize}
    \item Applied the Quantum Approximate Optimization Algorithm (QAOA) and quantum annealing to graph partitioning, to make high-precision cutting simulations faster and more accurate.
    \item On the first of two ANSYS blade models, the best run cut the worst-case number of solver iterations by about 64\%.
  \end{itemize}
  }

\cvrole
  {Research Assistant}
  {Mar 2024 -- Dec 2024}
  {Quantum Industrial Innovation, Senai Cimatec}
  {Salvador, Brazil}
  {
  \begin{itemize}
    \item Built the pipeline for the Facial Biosignals Time-Series Dataset (FBioT), which tracks facial feature movements for emotion recognition (\cvlink{https://doi.org/10.3390/electronics13244867}{paper}).
  \end{itemize}
  }

\cvrole
  {Research Assistant}
  {Mar 2024 -- Dec 2024}
  {Senai Cimatec}
  {Salvador, Brazil}
  {
  \begin{itemize}
    \item Modelled EEG time series with Detrended Fluctuation Analysis to identify the pre-ictal (pre-seizure) state in patients with epilepsy (\cvlink{https://orlandomotapires.github.io/assets/pdf/MISCOSSND.pdf}{paper}, \cvlink{https://github.com/orlandomotapires/Seizures-analysis}{code}).
  \end{itemize}
  }

\cvrole
  {Research Assistant}
  {Jan 2022 -- Dec 2023}
  {Senai Cimatec HPC Center}
  {Salvador, Brazil}
  {
  \begin{itemize}
    \item Tuned a CPU-based stencil computation with Data Parallel C++ (DPC++) on several devices (\cvlink{https://doi.org/10.34178/jbth.v6iSuppl1.269}{paper}).
    \item Co-wrote \cvlink{https://github.com/muriloboratto/AI-university-assignments}{introductory Python notebooks} on neural networks and artificial intelligence.
  \end{itemize}
  }

% --- SEÇÃO: HABILIDADES ---
\section{Skills}

\cvitem{Programming}{Python, C++, DPC++ / SYCL, Qiskit, PennyLane}
\cvitem{Quantum}{QAOA, quantum annealing, D-Wave, QLM}
\cvitem{Tools}{Git, Docker, Linux}

% --- SEÇÃO: PROJETOS ---
\clearpage
\section{Software}

\cvrole
  {HDD Studio}
  {2025 -- 2026}
  {Hybrid domain decomposition, the software behind my thesis}
  {}
  {
  \begin{itemize}
    \item A desktop app and command-line tool that reorders any square sparse matrix into diagonal blocks, by partitioning its graph with METIS, quantum annealing or QAOA, locally or on D-Wave and QLM quantum hardware (\cvlink{https://github.com/orlandomotapires/hybrid_domain_decomposition}{code}).
  \end{itemize}
  }

\cvrole
  {Ket Lab}
  {2026}
  {A quantum computing course that runs in the browser}
  {}
  {
  \begin{itemize}
    \item Write Qiskit or PennyLane code and watch the quantum state change line by line. 53 lessons in 7 tracks, from a single qubit to quantum annealing (\cvlink{https://orlandomotapires.github.io/ket-lab/}{open}, \cvlink{https://github.com/orlandomotapires/ket-lab}{code}).
  \end{itemize}
  }

\cvrole
  {CUDA Laplacian filter}
  {2023}
  {HPC hackathon, Fundamentals of Parallel Programming course}
  {}
  {
  \begin{itemize}
    \item Wrote the CUDA kernels for one and two GPUs (NCCL) and measured the run times. On a 46340 $\times$ 46340 input, the GPU version ran 19.68$\times$ faster than the sequential code (\cvlink{https://github.com/orlandomotapires/hackaton_fpp}{code}).
  \end{itemize}
  }

% --- SEÇÃO: ENSINO ---
\section{Teaching and Community}

\cvrole
  {Co-founder, Programming Club}
  {Sep 2021 -- Dec 2024}
  {Senai Cimatec}
  {Salvador, Brazil}
  {
  \begin{itemize}
    \item Taught algorithms and programming logic to engineering students, organised competitive programming marathons such as OBI and SBC, and ran the student initiative.
  \end{itemize}
  }

\cvrole
  {Intel Student Ambassador}
  {Mar 2023 -- Dec 2023}
  {Senai Cimatec}
  {Salvador, Brazil}
  {
  \begin{itemize}
    \item Ran workshops on Intel oneAPI for students at the university.
  \end{itemize}
  }

\cvrole
  {Teaching Assistant, Algorithms and Computational Thinking}
  {Feb 2022 -- Dec 2022}
  {Senai Cimatec}
  {Salvador, Brazil}
  {
  \begin{itemize}
    \item Graded assignments and exams, and taught extra classes.
  \end{itemize}
  }

% --- SEÇÃO: PUBLICAÇÕES ---
\section{Publications}
\small

\pub{Modeling Ictal Sentences Characterized by the Occurrence of Seizures in a Setting of Neurological Disorder}
{\textbf{Orlando Mota Pires}, Adrian Widmer, Florêncio Mendes Oliveira Filho}
{\textit{XXVIII Encontro Nacional de Modelagem Computacional (ENMC), 2025.} \cvlink{https://orlandomotapires.github.io/assets/pdf/MISCOSSND.pdf}{PDF}}

\pub{Facial Biosignals Time-Series Dataset (FBioT): A Visual-Temporal Facial Expression Recognition (VT-FER) Approach}
{João Marcelo Silva Souza, Caroline da Silva Morais Alves, Jés de Jesus Fiais Cerqueira, Wagner Luiz Alves de Oliveira, \textbf{Orlando Mota Pires} et al.}
{\textit{Electronics 13(24), 4867, 2024.} \cvlink{https://doi.org/10.3390/electronics13244867}{doi:10.3390/electronics13244867}}

\pub{Protótipo de Correlação de Eventos de Leitura Labial com Medidas Temporais e Legendas}
{\textbf{Orlando Mota Pires}, João Marcelo Silva Souza, Caroline da Silva M. Alves, Jés de Jesus F. Cerqueira, Wagner Luiz A. de Oliveira}
{\textit{XXV Congresso Brasileiro de Automática (CBA), 2024.}}

\pub{Tuning a CPU-Based Stencil Computation in a DPC++ Multi-Device Environment}
{Tiago Conceição Oliveira, Murilo Boratto, Antônio Horácio Rodrigues, \textbf{Orlando Mota Pires}, Leonardo Rodrigues Soares}
{\textit{Journal of Bioengineering, Technologies and Health 6(Suppl.\ 1), 24--30, 2023.} \cvlink{https://doi.org/10.34178/jbth.v6iSuppl1.269}{doi:10.34178/jbth.v6iSuppl1.269}}

\pub{OneAPI: an Approach for Developer-Centered Heterogeneous Computing}
{T. Conceição, A. H. Rodrigues, \textbf{O. Pires}, L. Soares}
{\textit{WSCAD, 2022.}}

\pub{Statistical Study of Eco-Efficiency in Compact and Average Cars (Chevrolet, Ford, VW, Fiat, Renault) in Brazil Based on the Metro Table in 2019}
{A. Wenceslau, A. Widmer, G. Dunkel, J. Samuel, \textbf{O. Pires}, A. Nascimento}
{\textit{Journal of Bioengineering, Technologies and Health, 2019.}}

\end{document}
